\chapter{Conclusion and Future Work}
\label{ch:conclusion}

\section{Conclusion}
% Write after the final runs. State the headline numbers for all three
% splits, the falling recall, whether the >=85% target was met on splits 1-2,
% and what the cross-subject gap was. Do not quote a single-split number.
[TODO]

\section{Limitations}

The following bound what the results in Chapter~\ref{ch:results} can be
taken to mean. They are properties of the design and the collection
conditions, and were known in advance rather than discovered afterwards.

\textbf{Single occupancy.} The system assumes one person in the sensing
area. With two or more, the channel perturbation is a superposition that a
classifier trained on single-occupancy data cannot decompose, and the
labelling problem itself becomes ambiguous.

\textbf{Amplitude only.} Raw CSI phase on a single-antenna ESP32-S3 is
corrupted by carrier and sampling frequency offsets, and the multi-antenna
differencing that would cancel them is not available on this hardware.
Phase-derived features --- including Doppler estimation, which is the most
physically direct route to a speed signal --- are therefore inaccessible.

\textbf{Environment sensitivity.} A model trained in one room learns that
room's multipath structure alongside the activity. The cross-environment
result quantifies this, and it is expected to be the weakest number in the
report. Redeployment into a new room should be assumed to require new data
until demonstrated otherwise.

\textbf{Dataset scale and subject diversity.} The dataset is collected by a
four-person team with recruited volunteers over the project's timeline. The
cross-subject result is therefore a leave-one-out estimate over a small
subject pool, whose variance across folds is itself substantial; it should
be read as an honest lower bound rather than a population-level claim.

\textbf{An uncontrolled radio environment.} The transmitter is an
ISP-supplied router whose channel cannot be pinned, operating on a 2.4\,GHz
channel shared with a neighbouring access point at near-equal power
(Table~\ref{tab:channel-survey}). Sessions are screened for channel drift,
but the contention and the uncontrolled multipath contributed by that
neighbour's traffic are present in the data and cannot be removed from it.

\textbf{Fixed geometry.} Receiver positions are measured, marked, and
photographed, and the model is trained on one geometry. Moving a receiver
invalidates the trained model in the same way that moving to a new room
does.

\textbf{Class boundary definition.} Transitions --- sitting down, standing
up, lying down --- are excluded rather than modelled, and windows straddling
a label boundary are discarded. A deployed system encounters transitions
continuously, so the reported accuracy is measured on cleaner label
conditions than continuous operation presents.

\section{Future Work}

\subsection{On-device inference}

Running quantized inference on the ESP32-S3 itself would remove the host
from the sensing loop entirely, turning each receiver into a self-contained
sensor. A budget analysis was carried out and establishes three things.

First, \textbf{weights are not the constraint}: the CNN quantizes to
approximately 236\,KB in int8, against 16\,MB of flash.

Second, \textbf{memory is constrained by the activation arena, not the
model}. The peak live working set for a 3\,s window is approximately
609\,KB, and it occurs entirely within the first convolutional block, whose
feature map is still at full time-subcarrier resolution. This exceeds the
S3's 512\,KB of internal SRAM and would have to be placed in the slower
8\,MB PSRAM. Because the peak is localised to the front end, a stride-2
first convolution would reduce both the arena and the arithmetic by roughly
a factor of four, bringing the model within internal SRAM --- at an accuracy
cost that would need to be measured and reported.

Third, and decisively, \textbf{on-device inference is necessarily
single-receiver}. A board holds only its own CSI; sharing between boards
would reintroduce the network dependency the change exists to remove. The
cost of the stretch goal is therefore not memory --- receivers are input
channels, and the parameter count barely moves --- but the loss of the
spatial diversity that motivated the three-receiver design. The
receiver-count ablation measures exactly this, and it is the number on which
the decision should turn.

\subsection{Other directions}

\textbf{Multi-person sensing} is the most requested extension and the
hardest, requiring either source separation or a fundamentally different
output formulation such as counting or coarse localisation.

\textbf{Through-wall operation} is established as feasible on this hardware
class \cite{strohmayer2023wall} and would extend coverage from one room to a
small dwelling.

\textbf{Domain adaptation} addresses the environment sensitivity above
directly: fine-tuning on a short unlabelled or lightly labelled capture in a
new room, rather than recollecting a full dataset, is what would make
deployment practical.

\textbf{Data augmentation} --- time warping, subcarrier dropout, amplitude
jitter --- is the natural first response to a cross-subject gap, and was
deliberately deferred until the size of that gap on real data was known, so
that the augmentation could be chosen against the observed failure rather
than assumed.

\textbf{Wireless transport} (UDP from receiver to host, instead of USB)
would remove the cabling that currently constrains receiver placement. The
distribution shift this introduces --- the receivers' own transmissions
become part of the channel --- would require a small adaptation set
collected in that mode.

\textbf{Releasing the dataset}, if the department permits, would make the
multi-receiver ESP32-S3 corpus a citable artefact and address the scarcity
of multi-receiver, multi-session CSI data in this hardware class.
